\lysh{32257}{4}
\begin{alist}
\item Η βαθμολογία των δέκα μαθητών σε αύξουσα σειρά είναι:
$7, 8, 9, 10, 10, 10, 11, 17, 18, 20$
Οπότε το εύρος της βαθμολογίας είναι: $R = x_{\max} - x_{\min} = 20 - 7 = 13.$
Εφόσον το πλήθος των παρατηρήσεων είναι 10, δηλαδή άρτιος αριθμός, η διάμεσος είναι το ημιάθροισμα των δύο μεσαίων παρατηρήσεων.
Δηλαδή $\delta = \frac{x_5 + x_6}{2} = \frac{10+10}{2} = 10.$
\item Για τη μέση τιμή της βαθμολογίας έχουμε:
\[ \bar{x} = \frac{\sum_{i=1}^{\nu} t_i}{\nu} = \frac{7+8+9+10+10+10+11+17+18+20}{10} = \frac{120}{10} = 12. \]

\item Η διακύμανση της βαθμολογίας δίνεται από τον τύπο: $s^2 = \frac{1}{\nu} \sum_{i=1}^{\nu} (x_i - \bar{x})^2$, από τον οποίο με αντικατάσταση παίρνουμε:
\[ \begin{aligned} s^2 &= \frac{1}{10} [(7-12)^2 + (8-12)^2 + (9-12)^2 + (10-12)^2 + (10-12)^2 + (10-12)^2 + (11-12)^2 \\ & \quad + (17-12)^2 + (18-12)^2 + (20-12)^2] \\ &= \frac{1}{10} (25+16+9+4+4+4+1+25+36+64) = \frac{1}{10} \cdot 188 = 18,8. \end{aligned} \]
Εφόσον η διακύμανση $s^2 = 18,8$ η τυπική απόκλιση $s = \sqrt{s^2} = \sqrt{18,8} \approx 4,3.$
Επομένως ο συντελεστής μεταβολής του δείγματος θα είναι:
\[ CV = \frac{s}{\bar{x}} = \frac{4,3}{12} \approx 0,36 = 36\% > 10\%, \]
οπότε το δείγμα δεν είναι ομοιογενές.

\item Γνωρίζουμε ότι αν σε κάθε βαθμό προσθέσουμε τον αριθμό 33, το νέο δείγμα αριθμών που θα προκύψει θα έχει μέση τιμή: $\bar{Y} = \bar{x}+33 = 12+33=45$ και τυπική απόκλιση $S_Y = s = 4,3.$
Επομένως ο συντελεστής μεταβολής του δείγματος θα είναι:
\[ CV_Y = \frac{S_Y}{\bar{Y}} = \frac{4,3}{45} \approx 0,095 = 9,5\% < 10\%, \]
οπότε το δείγμα των αριθμών που θα προκύψει θα είναι ομοιογενές. Άρα ο ισχυρισμός του μαθητή είναι σωστός.
\end{alist}
\vspace{6mm}
\noindent\rule{\textwidth}{0.4pt}
\vspace{6mm}

